\chapter{Lineare Regression}

\section{Ziel dieses Kapitels}

In diesem Kapitel betrachten wir lineare Regression (linear regression).
Lineare Regression ist eines der wichtigsten Grundmodelle im maschinellen Lernen.
Sie ist mathematisch überschaubar, aber viele zentrale Ideen späterer Modelle tauchen hier bereits auf:

\begin{itemize}
    \item Modellfunktion (model function),
    \item Parameter oder Gewichte (parameters, weights),
    \item Designmatrix (design matrix),
    \item Fehlerfunktion (error function),
    \item kleinste Quadrate (least squares),
    \item Maximum-Likelihood-Interpretation (maximum likelihood interpretation),
    \item Regularisierung (regularization),
    \item Bias-Term,
    \item Basisfunktionen (basis functions),
    \item Generalisierung (generalization).
\end{itemize}

\begin{conceptbox}
Lineare Regression (linear regression) ist ein überwachtes Lernverfahren (supervised learning) für kontinuierliche Zielwerte.

Das Modell sagt einen reellen Zielwert \(t\) aus einer Eingabe \(\vec{x}\) vorher.
\end{conceptbox}

Nach diesem Kapitel solltest du folgende Fragen beantworten können:

\begin{itemize}
    \item Was ist lineare Regression (linear regression)?
    \item Was bedeutet linear in den Parametern?
    \item Was ist ein Bias-Term?
    \item Was ist eine Designmatrix (design matrix)?
    \item Wie schreibt man lineare Regression in Matrixform?
    \item Was ist die Sum-of-Squares-Fehlerfunktion?
    \item Wie leitet man die Normalengleichung (normal equation) her?
    \item Wann ist die Lösung eindeutig?
    \item Wie hängt lineare Regression mit Maximum-Likelihood zusammen?
    \item Warum führt Gaußsches Rauschen zu quadratischem Fehler?
    \item Was ist Ridge Regression?
    \item Wie hängt Ridge Regression mit MAP zusammen?
    \item Was sind Basisfunktionen?
    \item Warum kann ein Modell nichtlinear in der Eingabe, aber linear in den Parametern sein?
\end{itemize}

\section{Regression als überwachtes Lernen}

Bei der Regression (regression) wollen wir einen kontinuierlichen Zielwert vorhersagen.
Ein Trainingsdatensatz besteht aus Eingaben und Zielwerten:

\[
    \mathcal{D}
    =
    \{(\vec{x}_1,t_1),(\vec{x}_2,t_2),\dots,(\vec{x}_N,t_N)\}.
\]

Dabei ist

\[
    \vec{x}_n \in \mathbb{R}^D
\]

der Eingabevektor des \(n\)-ten Datenpunkts und

\[
    t_n \in \mathbb{R}
\]

der zugehörige Zielwert.

Das Ziel ist, ein Modell \(y(\vec{x})\) zu lernen, sodass

\[
    y(\vec{x}_n) \approx t_n
\]

für die Trainingsdaten gilt und das Modell auch für neue Daten gut generalisiert.

\begin{definitionbox}
Regression (regression) ist ein überwachtes Lernproblem, bei dem der Zielwert kontinuierlich ist:

\[
    t \in \mathbb{R}.
\]
\end{definitionbox}

\begin{examplebox}
Beispiele für Regression:

\begin{itemize}
    \item Wohnungsdaten \(\vec{x}\) \(\rightarrow\) Mietpreis \(t\),
    \item Wetterdaten \(\vec{x}\) \(\rightarrow\) Temperatur \(t\),
    \item medizinische Messwerte \(\vec{x}\) \(\rightarrow\) Laborwert \(t\),
    \item Sensordaten \(\vec{x}\) \(\rightarrow\) erwartete Lebensdauer eines Bauteils \(t\).
\end{itemize}
\end{examplebox}

\section{Das einfachste lineare Modell}

Wir beginnen mit einer eindimensionalen Eingabe

\[
    x \in \mathbb{R}.
\]

Ein einfaches lineares Modell ist eine Gerade:

\[
    y(x,w_0,w_1)
    =
    w_0 + w_1x.
\]

Dabei sind \(w_0\) und \(w_1\) Parameter.

\begin{definitionbox}
Das Modell

\[
    y(x,w_0,w_1)
    =
    w_0 + w_1x
\]

heißt lineares Regressionsmodell für eine eindimensionale Eingabe.

Der Parameter \(w_0\) ist der Bias oder Achsenabschnitt (bias, intercept).
Der Parameter \(w_1\) ist die Steigung (slope).
\end{definitionbox}

\begin{conceptbox}
Der Bias-Term \(w_0\) erlaubt dem Modell, die Gerade nach oben oder unten zu verschieben.

Ohne Bias-Term hätte das Modell die Form

\[
    y(x,w_1)
    =
    w_1x
\]

und müsste immer durch den Ursprung gehen.
\end{conceptbox}

\section{Lineare Regression mit mehreren Eingabedimensionen}

In maschinellem Lernen ist die Eingabe meistens ein Vektor:

\[
    \vec{x}
    =
    \begin{pmatrix}
        x_1 \\
        x_2 \\
        \vdots \\
        x_D
    \end{pmatrix}
    \in \mathbb{R}^D.
\]

Ein lineares Modell mit Bias-Term ist:

\[
    y(\vec{x},\vec{w})
    =
    w_0
    +
    w_1x_1
    +
    w_2x_2
    +
    \dots
    +
    w_Dx_D.
\]

Hier sind

\[
    w_0,w_1,\dots,w_D
\]

die Parameter des Modells.

\begin{definitionbox}
Die lineare Regression für \(\vec{x}\in\mathbb{R}^D\) hat die Form

\[
    y(\vec{x},\vec{w})
    =
    w_0
    +
    \sum_{j=1}^{D}
    w_jx_j.
\]
\end{definitionbox}

Der Parameter \(w_0\) ist wieder der Bias-Term.
Die Parameter \(w_1,\dots,w_D\) bestimmen, wie stark die einzelnen Merkmale (features) auf die Vorhersage wirken.

\begin{examplebox}
Angenommen,

\[
    \vec{x}
    =
    \begin{pmatrix}
        \text{Wohnfläche} \\
        \text{Anzahl Zimmer} \\
        \text{Entfernung zum Zentrum}
    \end{pmatrix}.
\]

Dann könnte ein lineares Modell sein:

\[
    y(\vec{x},\vec{w})
    =
    w_0
    +
    w_1 \cdot \text{Wohnfläche}
    +
    w_2 \cdot \text{Anzahl Zimmer}
    +
    w_3 \cdot \text{Entfernung zum Zentrum}.
\]

Wenn \(w_1>0\), dann erhöht größere Wohnfläche die vorhergesagte Miete.
Wenn \(w_3<0\), dann senkt größere Entfernung zum Zentrum die vorhergesagte Miete.
\end{examplebox}

\section{Vektorform mit Bias}

Es ist praktisch, den Bias-Term in den Eingabevektor einzubauen.
Dazu erweitern wir den Eingabevektor um eine konstante Komponente:

\[
    \tilde{\vec{x}}
    =
    \begin{pmatrix}
        1 \\
        x_1 \\
        x_2 \\
        \vdots \\
        x_D
    \end{pmatrix}
    \in \mathbb{R}^{D+1}.
\]

Der Gewichtsvektor ist:

\[
    \vec{w}
    =
    \begin{pmatrix}
        w_0 \\
        w_1 \\
        w_2 \\
        \vdots \\
        w_D
    \end{pmatrix}
    \in \mathbb{R}^{D+1}.
\]

Dann kann das Modell kompakt geschrieben werden als:

\[
    y(\vec{x},\vec{w})
    =
    \vec{w}^{\top}\tilde{\vec{x}}.
\]

\begin{notebox}
Oft schreibt man wieder einfach \(\vec{x}\) statt \(\tilde{\vec{x}}\), obwohl die \(1\) für den Bias bereits eingefügt wurde.

Dann gilt:

\[
    y(\vec{x},\vec{w})
    =
    \vec{w}^{\top}\vec{x}.
\]

Dabei muss aus dem Kontext klar sein, ob \(\vec{x}\) den Bias-Eintrag enthält.
\end{notebox}

\section{Was bedeutet linear?}

Der Begriff linear in linearer Regression kann verwirrend sein.
Wichtig ist:

\begin{center}
\textbf{Lineare Regression ist linear in den Parametern, nicht notwendigerweise linear in der ursprünglichen Eingabe.}
\end{center}

Das einfache Modell

\[
    y(x,\vec{w})
    =
    w_0+w_1x
\]

ist linear in \(x\) und linear in den Parametern.

Aber auch das Modell

\[
    y(x,\vec{w})
    =
    w_0
    +
    w_1x
    +
    w_2x^2
\]

ist linear in den Parametern \(w_0,w_1,w_2\), obwohl es nichtlinear in \(x\) ist.

\begin{definitionbox}
Ein Modell ist linear in den Parametern, wenn es als Linearkombination von festen Funktionen geschrieben werden kann:

\[
    y(x,\vec{w})
    =
    \sum_{j=0}^{M-1}
    w_j \phi_j(x).
\]

Die Funktionen \(\phi_j(x)\) heißen Basisfunktionen (basis functions).
\end{definitionbox}

\begin{examplebox}
Für polynomiale Basisfunktionen gilt:

\[
    \phi_0(x)=1,
    \qquad
    \phi_1(x)=x,
    \qquad
    \phi_2(x)=x^2,
    \qquad
    \phi_3(x)=x^3.
\]

Dann ist

\[
    y(x,\vec{w})
    =
    w_0
    +
    w_1x
    +
    w_2x^2
    +
    w_3x^3.
\]

Dieses Modell ist linear in \(\vec{w}\), aber nichtlinear in \(x\).
\end{examplebox}

\section{Basisfunktionen}

Basisfunktionen (basis functions) transformieren Eingaben in neue Merkmale.
Statt direkt mit \(x\) zu arbeiten, arbeitet das Modell mit

\[
    \phi_0(x),\phi_1(x),\dots,\phi_{M-1}(x).
\]

Das Modell lautet:

\[
    y(x,\vec{w})
    =
    \sum_{j=0}^{M-1}
    w_j\phi_j(x).
\]

In Vektorform:

\[
    \vec{\phi}(x)
    =
    \begin{pmatrix}
        \phi_0(x) \\
        \phi_1(x) \\
        \vdots \\
        \phi_{M-1}(x)
    \end{pmatrix}.
\]

Dann:

\[
    y(x,\vec{w})
    =
    \vec{w}^{\top}\vec{\phi}(x).
\]

\begin{definitionbox}
Eine Basisfunktion (basis function) ist eine feste Funktion \(\phi_j(x)\), die aus der Eingabe ein neues Merkmal berechnet.

Das Modell bleibt linear in den Gewichten:

\[
    y(x,\vec{w})
    =
    \vec{w}^{\top}\vec{\phi}(x).
\]
\end{definitionbox}

\subsection{Typische Basisfunktionen}

Typische Basisfunktionen sind:

\begin{itemize}
    \item Polynomiale Basisfunktionen:
    \[
        \phi_j(x)=x^j.
    \]

    \item Gaußsche Basisfunktionen:
    \[
        \phi_j(x)
        =
        \exp\left(
            -\frac{(x-\mu_j)^2}{2s^2}
        \right).
    \]

    \item Sigmoid-Basisfunktionen:
    \[
        \phi_j(x)
        =
        \frac{1}{1+\exp(-a_j(x-b_j))}.
    \]
\end{itemize}

\begin{conceptbox}
Basisfunktionen erlauben nichtlineare Kurven, obwohl das Modell linear in den Parametern bleibt.

Dadurch können wir viele Methoden der linearen Algebra weiterverwenden.
\end{conceptbox}

\section{Trainingsdaten in Matrixform}

Wir betrachten nun \(N\) Trainingsdatenpunkte

\[
    (x_1,t_1),\dots,(x_N,t_N)
\]

und \(M\) Basisfunktionen.

Für jeden Datenpunkt \(x_n\) bilden wir den Feature-Vektor

\[
    \vec{\phi}(x_n)
    =
    \begin{pmatrix}
        \phi_0(x_n) \\
        \phi_1(x_n) \\
        \vdots \\
        \phi_{M-1}(x_n)
    \end{pmatrix}.
\]

Die Designmatrix (design matrix) \(\Phi\) enthält diese Feature-Vektoren als Zeilen:

\[
    \Phi
    =
    \begin{pmatrix}
        \phi_0(x_1) & \phi_1(x_1) & \dots & \phi_{M-1}(x_1) \\
        \phi_0(x_2) & \phi_1(x_2) & \dots & \phi_{M-1}(x_2) \\
        \vdots & \vdots & \ddots & \vdots \\
        \phi_0(x_N) & \phi_1(x_N) & \dots & \phi_{M-1}(x_N)
    \end{pmatrix}.
\]

Damit gilt:

\[
    \Phi \in \mathbb{R}^{N \times M}.
\]

Der Zielwertvektor ist:

\[
    \vec{t}
    =
    \begin{pmatrix}
        t_1 \\
        t_2 \\
        \vdots \\
        t_N
    \end{pmatrix}
    \in \mathbb{R}^{N}.
\]

Der Gewichtsvektor ist:

\[
    \vec{w}
    =
    \begin{pmatrix}
        w_0 \\
        w_1 \\
        \vdots \\
        w_{M-1}
    \end{pmatrix}
    \in \mathbb{R}^{M}.
\]

Die Vorhersagen für alle Datenpunkte sind:

\[
    \vec{y}
    =
    \Phi\vec{w}.
\]

\begin{definitionbox}
Die Designmatrix (design matrix) \(\Phi\) ist die Matrix aller transformierten Eingaben:

\[
    \Phi_{nj}
    =
    \phi_j(x_n).
\]

Jede Zeile entspricht einem Datenpunkt.
Jede Spalte entspricht einer Basisfunktion.
\end{definitionbox}

\begin{notebox}
Die Dimensionen sind wichtig:

\[
    \Phi \in \mathbb{R}^{N \times M},
    \qquad
    \vec{w} \in \mathbb{R}^{M},
    \qquad
    \vec{t} \in \mathbb{R}^{N}.
\]

Dann gilt:

\[
    \Phi\vec{w} \in \mathbb{R}^{N}.
\]
\end{notebox}

\section{Fehlerfunktion der linearen Regression}

Wir möchten, dass die Vorhersagen

\[
    y(x_n,\vec{w})
\]

möglichst nahe bei den Zielwerten \(t_n\) liegen.
Ein natürlicher Fehler ist die Differenz

\[
    y(x_n,\vec{w}) - t_n.
\]

Da positive und negative Fehler sich nicht gegenseitig aufheben sollen, quadrieren wir sie.
Die Sum-of-Squares-Fehlerfunktion ist:

\[
    E_D(\vec{w})
    =
    \frac{1}{2}
    \sum_{n=1}^{N}
    \left(
        y(x_n,\vec{w}) - t_n
    \right)^2.
\]

In Matrixform:

\[
    E_D(\vec{w})
    =
    \frac{1}{2}
    \left\|
        \Phi\vec{w}-\vec{t}
    \right\|^2.
\]

Das bedeutet:

\[
    E_D(\vec{w})
    =
    \frac{1}{2}
    (\Phi\vec{w}-\vec{t})^{\top}
    (\Phi\vec{w}-\vec{t}).
\]

\begin{definitionbox}
Die quadratische Fehlerfunktion (sum-of-squares error function) der linearen Regression ist

\[
    E_D(\vec{w})
    =
    \frac{1}{2}
    \sum_{n=1}^{N}
    \left(
        \vec{w}^{\top}\vec{\phi}(x_n)-t_n
    \right)^2.
\]

In Matrixform:

\[
    E_D(\vec{w})
    =
    \frac{1}{2}
    \left\|
        \Phi\vec{w}-\vec{t}
    \right\|^2.
\]
\end{definitionbox}

\begin{notebox}
Der Faktor \(\frac{1}{2}\) ändert die Lage des Minimums nicht.
Er ist nur praktisch, weil beim Ableiten des Quadrats ein Faktor \(2\) entsteht, der sich mit \(\frac{1}{2}\) kürzt.
\end{notebox}

\section{Least Squares}

Das Trainingsproblem der linearen Regression lautet:

\[
    \vec{w}^{\,*}
    =
    \arg\min_{\vec{w}}
    E_D(\vec{w}).
\]

Also:

\[
    \vec{w}^{\,*}
    =
    \arg\min_{\vec{w}}
    \frac{1}{2}
    \left\|
        \Phi\vec{w}-\vec{t}
    \right\|^2.
\]

\begin{definitionbox}
Least Squares oder kleinste Quadrate bezeichnet das Minimieren der Summe quadratischer Fehler:

\[
    \min_{\vec{w}}
    \sum_{n=1}^{N}
    \left(
        y(x_n,\vec{w})-t_n
    \right)^2.
\]
\end{definitionbox}

\section{Gradient der Fehlerfunktion}

Um das Minimum zu finden, leiten wir \(E_D(\vec{w})\) nach \(\vec{w}\) ab.

Wir starten mit:

\[
    E_D(\vec{w})
    =
    \frac{1}{2}
    (\Phi\vec{w}-\vec{t})^{\top}
    (\Phi\vec{w}-\vec{t}).
\]

Ausmultiplizieren:

\[
    E_D(\vec{w})
    =
    \frac{1}{2}
    \left(
        \vec{w}^{\top}\Phi^{\top}\Phi\vec{w}
        -
        \vec{w}^{\top}\Phi^{\top}\vec{t}
        -
        \vec{t}^{\top}\Phi\vec{w}
        +
        \vec{t}^{\top}\vec{t}
    \right).
\]

Da

\[
    \vec{w}^{\top}\Phi^{\top}\vec{t}
\]

ein Skalar ist, gilt:

\[
    \vec{w}^{\top}\Phi^{\top}\vec{t}
    =
    \vec{t}^{\top}\Phi\vec{w}.
\]

Also:

\[
    E_D(\vec{w})
    =
    \frac{1}{2}
    \left(
        \vec{w}^{\top}\Phi^{\top}\Phi\vec{w}
        -
        2\vec{w}^{\top}\Phi^{\top}\vec{t}
        +
        \vec{t}^{\top}\vec{t}
    \right).
\]

Wir verwenden die Ableitungsregeln:

\[
    \nabla_{\vec{w}}
    \left(
        \frac{1}{2}
        \vec{w}^{\top}A\vec{w}
    \right)
    =
    A\vec{w}
\]

für symmetrisches \(A\), und

\[
    \nabla_{\vec{w}}
    \left(
        \vec{w}^{\top}\vec{b}
    \right)
    =
    \vec{b}.
\]

Hier ist

\[
    A = \Phi^{\top}\Phi
\]

symmetrisch.
Damit:

\[
    \nabla_{\vec{w}}E_D(\vec{w})
    =
    \Phi^{\top}\Phi\vec{w}
    -
    \Phi^{\top}\vec{t}.
\]

\begin{definitionbox}
Der Gradient der quadratischen Fehlerfunktion ist

\[
    \nabla_{\vec{w}}E_D(\vec{w})
    =
    \Phi^{\top}
    (
        \Phi\vec{w}-\vec{t}
    ).
\]

Äquivalent:

\[
    \nabla_{\vec{w}}E_D(\vec{w})
    =
    \Phi^{\top}\Phi\vec{w}
    -
    \Phi^{\top}\vec{t}.
\]
\end{definitionbox}

\section{Normalengleichung}

Am Minimum muss der Gradient verschwinden:

\[
    \nabla_{\vec{w}}E_D(\vec{w})
    =
    \vec{0}.
\]

Also:

\[
    \Phi^{\top}\Phi\vec{w}
    -
    \Phi^{\top}\vec{t}
    =
    \vec{0}.
\]

Damit:

\[
    \Phi^{\top}\Phi\vec{w}
    =
    \Phi^{\top}\vec{t}.
\]

Diese Gleichung heißt Normalengleichung (normal equation).

\begin{definitionbox}
Die Normalengleichung der linearen Regression lautet:

\[
    \Phi^{\top}\Phi\vec{w}
    =
    \Phi^{\top}\vec{t}.
\]
\end{definitionbox}

Wenn \(\Phi^{\top}\Phi\) invertierbar ist, folgt:

\[
    \vec{w}_{\mathrm{LS}}
    =
    (\Phi^{\top}\Phi)^{-1}
    \Phi^{\top}\vec{t}.
\]

\begin{conceptbox}
Die Least-Squares-Lösung ist:

\[
    \vec{w}_{\mathrm{LS}}
    =
    (\Phi^{\top}\Phi)^{-1}
    \Phi^{\top}\vec{t}.
\]

Diese Formel gilt, wenn \(\Phi^{\top}\Phi\) invertierbar ist.
\end{conceptbox}

\section{Geometrische Interpretation}

Die Vorhersage

\[
    \vec{y}
    =
    \Phi\vec{w}
\]

liegt im Spaltenraum der Matrix \(\Phi\).
Der Zielvektor \(\vec{t}\) liegt im Allgemeinen nicht exakt in diesem Spaltenraum.

Least Squares sucht den Vektor \(\vec{y}\) im Spaltenraum von \(\Phi\), der \(\vec{t}\) am nächsten liegt.

Der Residuenvektor ist:

\[
    \vec{r}
    =
    \vec{t}
    -
    \Phi\vec{w}.
\]

Am Optimum ist der Residuenvektor orthogonal zum Spaltenraum von \(\Phi\):

\[
    \Phi^{\top}\vec{r}
    =
    \vec{0}.
\]

Einsetzen von \(\vec{r}=\vec{t}-\Phi\vec{w}\):

\[
    \Phi^{\top}
    (
        \vec{t}-\Phi\vec{w}
    )
    =
    \vec{0}.
\]

Das ist äquivalent zu:

\[
    \Phi^{\top}\Phi\vec{w}
    =
    \Phi^{\top}\vec{t}.
\]

\begin{notebox}
Die Normalengleichung bedeutet geometrisch:
Der Fehlervektor steht senkrecht auf allen Richtungen, die das Modell durch Änderung der Gewichte erreichen kann.
\end{notebox}

\section{Pseudoinverse}

Wenn \(\Phi^{\top}\Phi\) nicht invertierbar ist, kann man die Moore-Penrose-Pseudoinverse (Moore-Penrose pseudoinverse) verwenden.

Man schreibt:

\[
    \vec{w}_{\mathrm{LS}}
    =
    \Phi^{\dagger}\vec{t}.
\]

Falls \(\Phi^{\top}\Phi\) invertierbar ist, gilt:

\[
    \Phi^{\dagger}
    =
    (\Phi^{\top}\Phi)^{-1}\Phi^{\top}.
\]

\begin{definitionbox}
Die Pseudoinverse \(\Phi^{\dagger}\) verallgemeinert die Inverse auf nicht-quadratische oder singuläre Matrizen.

Die Least-Squares-Lösung kann geschrieben werden als

\[
    \vec{w}_{\mathrm{LS}}
    =
    \Phi^{\dagger}\vec{t}.
\]
\end{definitionbox}

\section{Wann ist die Lösung eindeutig?}

Die Lösung

\[
    \vec{w}_{\mathrm{LS}}
    =
    (\Phi^{\top}\Phi)^{-1}
    \Phi^{\top}\vec{t}
\]

ist eindeutig, wenn \(\Phi^{\top}\Phi\) invertierbar ist.
Das ist genau dann der Fall, wenn die Spalten von \(\Phi\) linear unabhängig sind.

\begin{conceptbox}
Für eine eindeutige Least-Squares-Lösung braucht man:

\[
    \mathrm{rank}(\Phi) = M.
\]

Das bedeutet:
Die \(M\) Spalten der Designmatrix müssen linear unabhängig sein.
\end{conceptbox}

Wenn \(M>N\), also mehr Parameter als Datenpunkte vorhanden sind, kann \(\Phi\) keinen vollen Spaltenrang haben.
Dann ist die Lösung nicht eindeutig.
Das ist ein typischer Fall von Unterbestimmtheit (underdetermined system).

\section{Training Error und Test Error}

Der Trainingsfehler (training error) ist der Fehler auf den Daten, die zum Lernen verwendet wurden:

\[
    E_{\mathrm{train}}
    =
    \frac{1}{2}
    \sum_{n=1}^{N}
    \left(
        y(\vec{x}_n,\vec{w})-t_n
    \right)^2.
\]

Der Testfehler (test error) wird auf neuen Daten gemessen:

\[
    E_{\mathrm{test}}
    =
    \frac{1}{2}
    \sum_{m=1}^{N_{\mathrm{test}}}
    \left(
        y(\vec{x}^{\,\mathrm{test}}_m,\vec{w})
        -
        t^{\mathrm{test}}_m
    \right)^2.
\]

\begin{conceptbox}
Ein kleiner Trainingsfehler bedeutet nicht automatisch gute Generalisierung.

Ein Modell kann die Trainingsdaten sehr gut treffen und trotzdem auf neuen Daten schlecht sein.
Das nennt man Überanpassung (overfitting).
\end{conceptbox}

\section{Polynomiale Regression und Overfitting}

Ein klassisches Beispiel für Overfitting ist polynomiale Regression.

Wir verwenden Basisfunktionen:

\[
    \phi_j(x)=x^j
\]

für

\[
    j=0,1,\dots,M-1.
\]

Das Modell ist:

\[
    y(x,\vec{w})
    =
    \sum_{j=0}^{M-1}
    w_jx^j.
\]

Wenn \(M\) klein ist, ist das Modell sehr eingeschränkt.
Es kann unteranpassen (underfitting).

Wenn \(M\) sehr groß ist, kann das Modell sehr flexibel werden.
Es kann dann die Trainingspunkte sehr genau treffen, aber zwischen den Punkten stark oszillieren.

\begin{conceptbox}
Kleines \(M\):

\[
    \text{einfaches Modell}
    \quad
    \Rightarrow
    \quad
    \text{mögliche Unteranpassung}.
\]

Großes \(M\):

\[
    \text{flexibles Modell}
    \quad
    \Rightarrow
    \quad
    \text{mögliche Überanpassung}.
\]
\end{conceptbox}

\section{Probabilistische Interpretation}

Lineare Regression kann probabilistisch interpretiert werden.
Wir nehmen an, dass die Zielwerte aus einer Modellvorhersage plus Rauschen entstehen:

\[
    t_n
    =
    y(\vec{x}_n,\vec{w})
    +
    \epsilon_n.
\]

Das Rauschen sei normalverteilt:

\[
    \epsilon_n
    \sim
    \mathcal{N}(0,\beta^{-1}).
\]

Dabei ist \(\beta\) die Präzision (precision), also der Kehrwert der Varianz:

\[
    \beta
    =
    \frac{1}{\sigma^2}.
\]

Dann gilt:

\[
    p(t_n \mid \vec{x}_n,\vec{w},\beta)
    =
    \mathcal{N}
    \left(
        t_n
        \mid
        y(\vec{x}_n,\vec{w}),
        \beta^{-1}
    \right).
\]

\begin{definitionbox}
Das probabilistische Modell der linearen Regression ist:

\[
    p(t \mid \vec{x},\vec{w},\beta)
    =
    \mathcal{N}
    \left(
        t
        \mid
        \vec{w}^{\top}\vec{\phi}(\vec{x}),
        \beta^{-1}
    \right).
\]
\end{definitionbox}

\section{Likelihood der linearen Regression}

Für i.i.d. Daten gilt:

\[
    p(\vec{t} \mid X,\vec{w},\beta)
    =
    \prod_{n=1}^{N}
    \mathcal{N}
    \left(
        t_n
        \mid
        \vec{w}^{\top}\vec{\phi}(\vec{x}_n),
        \beta^{-1}
    \right).
\]

Die Normalverteilungsdichte ist:

\[
    \mathcal{N}
    \left(
        t_n
        \mid
        y_n,
        \beta^{-1}
    \right)
    =
    \left(
        \frac{\beta}{2\pi}
    \right)^{1/2}
    \exp\left(
        -\frac{\beta}{2}
        (t_n-y_n)^2
    \right),
\]

wobei

\[
    y_n
    =
    \vec{w}^{\top}\vec{\phi}(\vec{x}_n).
\]

Damit:

\[
    p(\vec{t} \mid X,\vec{w},\beta)
    =
    \prod_{n=1}^{N}
    \left(
        \frac{\beta}{2\pi}
    \right)^{1/2}
    \exp\left(
        -\frac{\beta}{2}
        (t_n-y_n)^2
    \right).
\]

\section{Log-Likelihood der linearen Regression}

Die Log-Likelihood ist:

\[
    \log p(\vec{t} \mid X,\vec{w},\beta)
    =
    \sum_{n=1}^{N}
    \left[
        \frac{1}{2}\log\beta
        -
        \frac{1}{2}\log(2\pi)
        -
        \frac{\beta}{2}
        (t_n-y_n)^2
    \right].
\]

Zusammenfassen:

\[
    \log p(\vec{t} \mid X,\vec{w},\beta)
    =
    \frac{N}{2}\log\beta
    -
    \frac{N}{2}\log(2\pi)
    -
    \frac{\beta}{2}
    \sum_{n=1}^{N}
    (t_n-y_n)^2.
\]

Für festes \(\beta\) hängen die ersten beiden Terme nicht von \(\vec{w}\) ab.
Also maximiert man die Log-Likelihood nach \(\vec{w}\), indem man minimiert:

\[
    \sum_{n=1}^{N}
    (t_n-y_n)^2.
\]

\begin{conceptbox}
Maximum-Likelihood für lineare Regression mit Gaußschem Rauschen ist äquivalent zu Least Squares.

\[
    \text{Gauß-Likelihood}
    \quad
    \Longleftrightarrow
    \quad
    \text{quadratischer Fehler}.
\]
\end{conceptbox}

\section{Schätzung der Rauschvarianz}

Neben \(\vec{w}\) kann auch \(\beta\) geschätzt werden.
Die Maximum-Likelihood-Lösung für die Varianz ist:

\[
    \hat{\sigma}^2_{\mathrm{ML}}
    =
    \frac{1}{N}
    \sum_{n=1}^{N}
    \left(
        t_n-y(\vec{x}_n,\vec{w}_{\mathrm{ML}})
    \right)^2.
\]

Da

\[
    \beta = \frac{1}{\sigma^2},
\]

gilt:

\[
    \hat{\beta}_{\mathrm{ML}}
    =
    \frac{1}{\hat{\sigma}^2_{\mathrm{ML}}}.
\]

Also:

\[
    \hat{\beta}_{\mathrm{ML}}^{-1}
    =
    \frac{1}{N}
    \sum_{n=1}^{N}
    \left(
        t_n-y(\vec{x}_n,\vec{w}_{\mathrm{ML}})
    \right)^2.
\]

\begin{notebox}
Wie bei der Normalverteilung ist der ML-Schätzer der Varianz bei endlichem \(N\) verzerrt.
Für reine Modellschätzung wird trotzdem oft die ML-Form verwendet.
\end{notebox}

\section{Regularisierung}

Bei flexiblen Modellen kann Least Squares zu Overfitting führen.
Eine typische Lösung ist Regularisierung (regularization).

Dabei ergänzt man die Fehlerfunktion um einen Term, der große Gewichte bestraft.

\[
    E(\vec{w})
    =
    E_D(\vec{w})
    +
    \lambda E_W(\vec{w}).
\]

Hier ist:

\[
    E_D(\vec{w})
\]

der Datenfehler (data error), und

\[
    E_W(\vec{w})
\]

der Regularisierungsterm (regularization term).

\begin{conceptbox}
Regularisierung verhindert, dass die Gewichte zu groß werden.

Große Gewichte können dazu führen, dass das Modell stark schwankt und die Trainingsdaten überanpasst.
\end{conceptbox}

\section{Ridge Regression}

Eine häufige Regularisierung ist \(L_2\)-Regularisierung:

\[
    E_W(\vec{w})
    =
    \frac{1}{2}
    \vec{w}^{\top}\vec{w}.
\]

Dann ist die gesamte Fehlerfunktion:

\[
    E(\vec{w})
    =
    \frac{1}{2}
    \left\|
        \Phi\vec{w}-\vec{t}
    \right\|^2
    +
    \frac{\lambda}{2}
    \vec{w}^{\top}\vec{w}.
\]

Das nennt man Ridge Regression.

\begin{definitionbox}
Ridge Regression minimiert:

\[
    E(\vec{w})
    =
    \frac{1}{2}
    \left\|
        \Phi\vec{w}-\vec{t}
    \right\|^2
    +
    \frac{\lambda}{2}
    \vec{w}^{\top}\vec{w}.
\]

Der Parameter \(\lambda \geq 0\) steuert die Stärke der Regularisierung.
\end{definitionbox}

\section{Ridge-Lösung}

Wir leiten die Lösung her.
Die Fehlerfunktion ist:

\[
    E(\vec{w})
    =
    \frac{1}{2}
    (\Phi\vec{w}-\vec{t})^{\top}
    (\Phi\vec{w}-\vec{t})
    +
    \frac{\lambda}{2}
    \vec{w}^{\top}\vec{w}.
\]

Der Gradient ist:

\[
    \nabla_{\vec{w}}E(\vec{w})
    =
    \Phi^{\top}\Phi\vec{w}
    -
    \Phi^{\top}\vec{t}
    +
    \lambda\vec{w}.
\]

Setze den Gradienten gleich \(0\):

\[
    \Phi^{\top}\Phi\vec{w}
    -
    \Phi^{\top}\vec{t}
    +
    \lambda\vec{w}
    =
    \vec{0}.
\]

Zusammenfassen:

\[
    (\Phi^{\top}\Phi+\lambda I)\vec{w}
    =
    \Phi^{\top}\vec{t}.
\]

Wenn \(\Phi^{\top}\Phi+\lambda I\) invertierbar ist, folgt:

\[
    \vec{w}_{\mathrm{ridge}}
    =
    (\Phi^{\top}\Phi+\lambda I)^{-1}
    \Phi^{\top}\vec{t}.
\]

\begin{conceptbox}
Die Ridge-Lösung ist:

\[
    \vec{w}_{\mathrm{ridge}}
    =
    (\Phi^{\top}\Phi+\lambda I)^{-1}
    \Phi^{\top}\vec{t}.
\]

Für \(\lambda>0\) ist die Matrix oft besser konditioniert als \(\Phi^{\top}\Phi\).
\end{conceptbox}

\begin{notebox}
In manchen Konventionen wird der Bias-Term \(w_0\) nicht regularisiert.
Dann ersetzt man \(\lambda I\) durch eine Diagonalmatrix, deren erster Eintrag \(0\) ist.
\end{notebox}

\section{MAP-Interpretation von Ridge Regression}

Ridge Regression kann probabilistisch als MAP-Schätzung interpretiert werden.

Wir verwenden die Likelihood:

\[
    p(\vec{t}\mid X,\vec{w},\beta)
    =
    \prod_{n=1}^{N}
    \mathcal{N}
    \left(
        t_n
        \mid
        \vec{w}^{\top}\vec{\phi}(\vec{x}_n),
        \beta^{-1}
    \right).
\]

Außerdem verwenden wir einen Gauß-Prior auf die Gewichte:

\[
    p(\vec{w}\mid\alpha)
    =
    \mathcal{N}
    \left(
        \vec{w}
        \mid
        \vec{0},
        \alpha^{-1}I
    \right).
\]

Der negative Log-Posterior ist bis auf Konstanten:

\[
    \frac{\beta}{2}
    \left\|
        \Phi\vec{w}-\vec{t}
    \right\|^2
    +
    \frac{\alpha}{2}
    \vec{w}^{\top}\vec{w}.
\]

Multipliziert man mit \(\frac{1}{\beta}\), erhält man ein äquivalentes Optimierungsproblem:

\[
    \frac{1}{2}
    \left\|
        \Phi\vec{w}-\vec{t}
    \right\|^2
    +
    \frac{\alpha}{2\beta}
    \vec{w}^{\top}\vec{w}.
\]

Also entspricht Ridge Regression einem Gauß-Prior mit

\[
    \lambda
    =
    \frac{\alpha}{\beta}.
\]

\begin{conceptbox}
Ridge Regression ist MAP-Schätzung bei:

\[
    \text{Gauß-Likelihood}
    +
    \text{Gauß-Prior auf } \vec{w}.
\]

Der Regularisierungsparameter ist:

\[
    \lambda
    =
    \frac{\alpha}{\beta}.
\]
\end{conceptbox}

\section{Gradientenabstieg für lineare Regression}

Obwohl lineare Regression eine geschlossene Lösung hat, kann man sie auch iterativ mit Gradientenabstieg (gradient descent) trainieren.

Die Fehlerfunktion ist:

\[
    E_D(\vec{w})
    =
    \frac{1}{2}
    \left\|
        \Phi\vec{w}-\vec{t}
    \right\|^2.
\]

Der Gradient ist:

\[
    \nabla_{\vec{w}}E_D(\vec{w})
    =
    \Phi^{\top}
    (
        \Phi\vec{w}-\vec{t}
    ).
\]

Gradientenabstieg aktualisiert:

\[
    \vec{w}^{(\tau+1)}
    =
    \vec{w}^{(\tau)}
    -
    \eta
    \nabla_{\vec{w}}E_D(\vec{w}^{(\tau)}),
\]

also:

\[
    \vec{w}^{(\tau+1)}
    =
    \vec{w}^{(\tau)}
    -
    \eta
    \Phi^{\top}
    (
        \Phi\vec{w}^{(\tau)}-\vec{t}
    ).
\]

Dabei ist \(\eta>0\) die Lernrate (learning rate).

\begin{notebox}
Für große Datensätze ist eine iterative Optimierung oft praktischer als das direkte Invertieren von \(\Phi^{\top}\Phi\).

Das ist besonders wichtig bei sehr vielen Datenpunkten oder sehr vielen Parametern.
\end{notebox}

\section{Stochastischer Gradientenabstieg}

Beim stochastischen Gradientenabstieg (stochastic gradient descent, SGD) verwendet man nicht alle Daten auf einmal.
Stattdessen aktualisiert man die Gewichte mit einem einzelnen Datenpunkt oder einem Mini-Batch.

Für einen einzelnen Datenpunkt ist der Fehler:

\[
    E_n(\vec{w})
    =
    \frac{1}{2}
    \left(
        y(\vec{x}_n,\vec{w})-t_n
    \right)^2.
\]

Der Gradient ist:

\[
    \nabla_{\vec{w}}E_n(\vec{w})
    =
    \left(
        y(\vec{x}_n,\vec{w})-t_n
    \right)
    \vec{\phi}(\vec{x}_n).
\]

Das Update ist:

\[
    \vec{w}^{(\tau+1)}
    =
    \vec{w}^{(\tau)}
    -
    \eta
    \left(
        y(\vec{x}_n,\vec{w}^{(\tau)})-t_n
    \right)
    \vec{\phi}(\vec{x}_n).
\]

\begin{conceptbox}
SGD ist wichtig, weil es auf sehr große Datensätze skaliert.
Viele moderne ML-Modelle, insbesondere neuronale Netze, werden mit Varianten von SGD trainiert.
\end{conceptbox}

\section{Mehrdimensionale Zielwerte}

Bisher war der Zielwert eindimensional:

\[
    t_n \in \mathbb{R}.
\]

Man kann lineare Regression auch für mehrere Zielwerte verwenden:

\[
    \vec{t}_n \in \mathbb{R}^K.
\]

Dann verwendet man eine Gewichtsmatrix \(W\):

\[
    \vec{y}(\vec{x})
    =
    W^{\top}\vec{\phi}(\vec{x}).
\]

Die Zielmatrix ist:

\[
    T
    =
    \begin{pmatrix}
        \vec{t}_1^{\top} \\
        \vec{t}_2^{\top} \\
        \vdots \\
        \vec{t}_N^{\top}
    \end{pmatrix}
    \in \mathbb{R}^{N \times K}.
\]

Die Vorhersagematrix ist:

\[
    Y
    =
    \Phi W.
\]

Die Least-Squares-Lösung ist:

\[
    W_{\mathrm{LS}}
    =
    (\Phi^{\top}\Phi)^{-1}
    \Phi^{\top}T.
\]

\begin{notebox}
Mehrdimensionale lineare Regression kann als mehrere lineare Regressionsprobleme gleichzeitig verstanden werden.
Jede Ausgabedimension hat ihre eigenen Gewichte.
\end{notebox}

\section{Praktische Aspekte}

\subsection{Skalierung der Eingaben}

Lineare Regression ist empfindlich gegenüber Skalierung.
Wenn ein Merkmal sehr große Werte hat und ein anderes sehr kleine, können numerische Probleme entstehen.

Eine typische Standardisierung ist:

\[
    x_j'
    =
    \frac{x_j-\mu_j}{\sigma_j}.
\]

Dabei sind \(\mu_j\) und \(\sigma_j\) Mittelwert und Standardabweichung des \(j\)-ten Merkmals im Trainingsdatensatz.

\begin{examtipbox}
Wichtig:
Die Skalierungsparameter \(\mu_j\) und \(\sigma_j\) werden nur auf den Trainingsdaten berechnet.
Danach werden dieselben Werte auf Validierungs- und Testdaten angewendet.
\end{examtipbox}

\subsection{Multikollinearität}

Wenn Spalten der Designmatrix stark linear abhängig sind, ist \(\Phi^{\top}\Phi\) schlecht konditioniert.
Dann können kleine Änderungen in den Daten große Änderungen in den Gewichten verursachen.

Ridge Regression hilft hier, weil

\[
    \Phi^{\top}\Phi+\lambda I
\]

stabiler invertierbar ist.

\section{Häufige Fehler}

\begin{examtipbox}
Typische Fehler bei linearer Regression:

\begin{enumerate}
    \item Den Bias-Term \(w_0\) vergessen.
    \item Verwechseln, ob \(\vec{x}\) den Bias-Eintrag \(1\) enthält.
    \item Denken, dass lineare Regression immer linear in der ursprünglichen Eingabe sein muss.
    \item Designmatrix \(\Phi\) falsch dimensionieren.
    \item \(\Phi\vec{w}\) und \(\vec{w}^{\top}\Phi\) verwechseln.
    \item Die Normalengleichung falsch herum schreiben.
    \item Vergessen, dass \(\Phi^{\top}\Phi\) invertierbar sein muss.
    \item Overfitting bei zu vielen Basisfunktionen ignorieren.
    \item Den Zusammenhang zwischen Gaußschem Rauschen und quadratischem Fehler nicht erkennen.
    \item Ridge Regression nur als Trick sehen, nicht als MAP mit Gauß-Prior.
    \item Trainingsfehler mit Testfehler verwechseln.
\end{enumerate}
\end{examtipbox}

\section{Mini-Aufgaben}

\subsection{Aufgabe 1: Einfache Vorhersage}

Gegeben sei das Modell

\[
    y(x,w_0,w_1)
    =
    w_0+w_1x.
\]

Mit

\[
    w_0=2,
    \qquad
    w_1=3,
    \qquad
    x=4
\]

berechne \(y\).

\begin{examplebox}
Einsetzen:

\[
    y
    =
    2+3\cdot 4
    =
    2+12
    =
    14.
\]
\end{examplebox}

\subsection{Aufgabe 2: Designmatrix}

Gegeben seien die Datenpunkte

\[
    x_1=1,
    \qquad
    x_2=2,
    \qquad
    x_3=3.
\]

Wir verwenden Basisfunktionen

\[
    \phi_0(x)=1,
    \qquad
    \phi_1(x)=x.
\]

Schreibe die Designmatrix \(\Phi\).

\begin{examplebox}
Die Designmatrix hat eine Zeile pro Datenpunkt:

\[
    \Phi
    =
    \begin{pmatrix}
        1 & 1 \\
        1 & 2 \\
        1 & 3
    \end{pmatrix}.
\]
\end{examplebox}

\subsection{Aufgabe 3: Quadratischer Fehler}

Gegeben seien Zielwerte

\[
    \vec{t}
    =
    \begin{pmatrix}
        1 \\
        2
    \end{pmatrix}
\]

und Vorhersagen

\[
    \vec{y}
    =
    \begin{pmatrix}
        2 \\
        4
    \end{pmatrix}.
\]

Berechne

\[
    E_D
    =
    \frac{1}{2}
    \|\vec{y}-\vec{t}\|^2.
\]

\begin{examplebox}
Zuerst:

\[
    \vec{y}-\vec{t}
    =
    \begin{pmatrix}
        2 \\
        4
    \end{pmatrix}
    -
    \begin{pmatrix}
        1 \\
        2
    \end{pmatrix}
    =
    \begin{pmatrix}
        1 \\
        2
    \end{pmatrix}.
\]

Dann:

\[
    \|\vec{y}-\vec{t}\|^2
    =
    1^2+2^2
    =
    5.
\]

Also:

\[
    E_D
    =
    \frac{1}{2}
    5
    =
    2{,}5.
\]
\end{examplebox}

\subsection{Aufgabe 4: Ridge-Lösung}

Gegeben sei

\[
    \Phi^{\top}\Phi
    =
    \begin{pmatrix}
        2 & 0 \\
        0 & 3
    \end{pmatrix},
    \qquad
    \Phi^{\top}\vec{t}
    =
    \begin{pmatrix}
        4 \\
        9
    \end{pmatrix},
\]

und

\[
    \lambda = 1.
\]

Berechne \(\vec{w}_{\mathrm{ridge}}\).

\begin{examplebox}
Ridge Regression verwendet:

\[
    \vec{w}_{\mathrm{ridge}}
    =
    (\Phi^{\top}\Phi+\lambda I)^{-1}
    \Phi^{\top}\vec{t}.
\]

Zuerst:

\[
    \Phi^{\top}\Phi+\lambda I
    =
    \begin{pmatrix}
        2 & 0 \\
        0 & 3
    \end{pmatrix}
    +
    \begin{pmatrix}
        1 & 0 \\
        0 & 1
    \end{pmatrix}
    =
    \begin{pmatrix}
        3 & 0 \\
        0 & 4
    \end{pmatrix}.
\]

Die Inverse ist:

\[
    (\Phi^{\top}\Phi+\lambda I)^{-1}
    =
    \begin{pmatrix}
        \frac{1}{3} & 0 \\
        0 & \frac{1}{4}
    \end{pmatrix}.
\]

Also:

\[
    \vec{w}_{\mathrm{ridge}}
    =
    \begin{pmatrix}
        \frac{1}{3} & 0 \\
        0 & \frac{1}{4}
    \end{pmatrix}
    \begin{pmatrix}
        4 \\
        9
    \end{pmatrix}
    =
    \begin{pmatrix}
        \frac{4}{3} \\
        \frac{9}{4}
    \end{pmatrix}.
\]
\end{examplebox}

\section{Zusammenfassung}

\begin{itemize}
    \item Lineare Regression (linear regression) ist ein überwachtes Lernverfahren für kontinuierliche Zielwerte.
    \item Das einfache lineare Modell ist:
    \[
        y(x,w_0,w_1)
        =
        w_0+w_1x.
    \]
    \item Mit mehreren Merkmalen:
    \[
        y(\vec{x},\vec{w})
        =
        w_0+\sum_{j=1}^{D}w_jx_j.
    \]
    \item Mit Bias in Vektorform:
    \[
        y(\vec{x},\vec{w})
        =
        \vec{w}^{\top}\vec{x}.
    \]
    \item Lineare Regression ist linear in den Parametern, nicht notwendigerweise linear in der Eingabe.
    \item Mit Basisfunktionen:
    \[
        y(x,\vec{w})
        =
        \sum_{j=0}^{M-1}
        w_j\phi_j(x).
    \]
    \item Die Designmatrix ist:
    \[
        \Phi_{nj}
        =
        \phi_j(x_n).
    \]
    \item Die Vorhersagen für alle Datenpunkte sind:
    \[
        \vec{y}
        =
        \Phi\vec{w}.
    \]
    \item Die quadratische Fehlerfunktion ist:
    \[
        E_D(\vec{w})
        =
        \frac{1}{2}
        \|\Phi\vec{w}-\vec{t}\|^2.
    \]
    \item Die Normalengleichung lautet:
    \[
        \Phi^{\top}\Phi\vec{w}
        =
        \Phi^{\top}\vec{t}.
    \]
    \item Die Least-Squares-Lösung ist:
    \[
        \vec{w}_{\mathrm{LS}}
        =
        (\Phi^{\top}\Phi)^{-1}
        \Phi^{\top}\vec{t}.
    \]
    \item Mit Gaußschem Rauschen ist Least Squares äquivalent zu Maximum-Likelihood.
    \item Ridge Regression minimiert:
    \[
        E(\vec{w})
        =
        \frac{1}{2}
        \|\Phi\vec{w}-\vec{t}\|^2
        +
        \frac{\lambda}{2}
        \vec{w}^{\top}\vec{w}.
    \]
    \item Die Ridge-Lösung ist:
    \[
        \vec{w}_{\mathrm{ridge}}
        =
        (\Phi^{\top}\Phi+\lambda I)^{-1}
        \Phi^{\top}\vec{t}.
    \]
    \item Ridge Regression entspricht MAP mit Gauß-Likelihood und Gauß-Prior auf die Gewichte.
\end{itemize}

\section{Typische Prüfungsfragen}

\begin{examtipbox}
Mögliche Prüfungsfragen zu diesem Kapitel:

\begin{enumerate}
    \item Was ist lineare Regression?
    \item Was ist der Unterschied zwischen Klassifikation und Regression?
    \item Was ist ein Bias-Term?
    \item Warum erweitert man den Eingabevektor um eine \(1\)?
    \item Was bedeutet linear in den Parametern?
    \item Gib ein Beispiel für ein Modell, das nichtlinear in \(x\), aber linear in \(\vec{w}\) ist.
    \item Was ist eine Basisfunktion?
    \item Was ist die Designmatrix?
    \item Welche Dimensionen haben \(\Phi\), \(\vec{w}\) und \(\vec{t}\)?
    \item Schreibe die quadratische Fehlerfunktion in Summenform und Matrixform.
    \item Leite den Gradienten von \(E_D(\vec{w})\) her.
    \item Leite die Normalengleichung her.
    \item Wann ist die Least-Squares-Lösung eindeutig?
    \item Was ist die Pseudoinverse?
    \item Erkläre die geometrische Interpretation der Normalengleichung.
    \item Warum kann polynomiale Regression überanpassen?
    \item Wie hängt lineare Regression mit Maximum-Likelihood zusammen?
    \item Warum führt Gaußsches Rauschen zu quadratischem Fehler?
    \item Was ist Ridge Regression?
    \item Leite die Ridge-Lösung her.
    \item Wie hängt Ridge Regression mit MAP zusammen?
    \item Was ist der Unterschied zwischen Trainingsfehler und Testfehler?
\end{enumerate}
\end{examtipbox}